\documentclass[12pt]{article}
\usepackage[utf8]{inputenc}
\usepackage{amsmath, amssymb, amsthm}
\usepackage{geometry}
\usepackage{fancyhdr}

\geometry{letterpaper, margin=1in}
\setlength{\headheight}{15pt}

% Header setup
\pagestyle{fancy}
\fancyhf{}
\lhead{Math 477 Spring 2026}
\rhead{Jeffrey Lu}
\cfoot{\thepage}

\title{Homework Assignment 1 \\ \large Math 292, Spring 2026 \\ Rutgers University}
\author{Jeffrey Lu}

\begin{document}

\maketitle

\noindent \textbf{Problem 1:}
\begin{itemize}
    \item [(a)]
          \begin{align*}
              E[(3+2Y)^2] & = E[9 + 12Y + 4Y^2]              \\
                          & = 9 + 12E[Y] + 4E[Y^2]           \\
                          & = 9 + 12(2) + 4(Var(Y) + E[Y]^2) \\
                          & = 65
          \end{align*}

    \item [(b)]
          \begin{align*}
              Var(5+4Y) & = 4^2Var(Y) = 64
          \end{align*}
\end{itemize}

\noindent \textbf{Problem 2:}
\begin{itemize}
    \item [(a)]
          \begin{align*}
              P(X=1) = p(HT) + p(TH) = 0.6 * 0.3 + 0.4 * 0.7 = 0.46
          \end{align*}

    \item [(b)]
          \begin{align*}
              E[X] & = 0 * P(X=0) + 1 * P(X=1) + 2 * P(X=2) \\
                   & = 0.46 + 2 * 0.42                      \\
                   & = 1.3
          \end{align*}
\end{itemize}

\noindent \textbf{Problem 3:}

\begin{itemize}
    \item [(a)]
          \begin{align*}
              p(\text{Strategy 1}) & = p                               \\
              p(\text{Strategy 2}) & = p^2 + \frac{1}{2} * 2p(1-p) = p
          \end{align*}

    \item [(b)] Strategy 2 has a equal probability as Strategy 1.
    \item [(c)]
          \begin{align*}
              E[X] = 10 * 0.7 = 7
          \end{align*}
\end{itemize}

\noindent \textbf{Problem 4:}

\begin{itemize}
    \item [(a)]
          Let A be the event that dice A is rolled, and B the event that dice B is rolled and G be the event that green is rolled
          \begin{align*}
              P(G) & = P(G ~|~ A)P(A) + P(G ~|~ B)P(B)                       \\
                   & = \frac{1}{2} * \frac{1}{2} + \frac{1}{6} * \frac{1}{2} \\
                   & = \frac{1}{3}
          \end{align*}

    \item [(b)]

          The probability of rolling two greens in a row is
          \begin{align*}
              P(G_1, G_2) & = P(G_1, G_2 ~|~ A)P(A) + P(G_1, G_2 ~|~ B)P(B)  \\
                          & = (\frac{1}{2})^3 + (\frac{1}{6})^2(\frac{1}{2}) \\
                          & = \frac{5}{36}
          \end{align*}

          Then we have that
          \begin{align*}
              P(G_3 ~|~ G_1, G_2) & =  \frac{P(G_1, G_2,G_3 ~|~ A)P(A) + P(G_1, G_2, G_3 ~|~ B)P(B)}{ P(G_1, G_2)} \\
                                  & = \frac{(\frac{1}{2})^4 + (\frac{1}{6})^3(\frac{1}{2})}{\frac{5}{36}}          \\
                                  & = \frac{7}{15}
          \end{align*}
    \item [(c)]
          \begin{align*}
              P(A ~|~ G_1, G_2) & = \frac{P(G_1, G_2 ~|~ A) P(A)}{P(G_1, G_2)} \\
                                & = \frac{(\frac{1}{2})^3}{\frac{5}{36}}       \\
                                & = \frac{9}{10}
          \end{align*}

\end{itemize}

\noindent \textbf{Problem 5:}

Let W be the event of winning \$1.10 and it's complement losing \$1.00

Then, we can solve for the probability of two of the same color marbles being picked
\begin{align*}
    P(W)   & = \frac{\binom{5}{2} + \binom{5}{2}}{\binom{10}{2}} = \frac{4}{9} \\
    P(W^c) & = 1 - P(W) = \frac{5}{9}
\end{align*}

Therefore, the expected value is just
\begin{align*}
    E[X] = 1.1 * P(W) + -1 * P(W^c) \approx{-0.067}
\end{align*}

Then, the variance is
\begin{align*}
    Var[X] = E[X^2] - E[X]^2 \approx{1.0889}
\end{align*}

\noindent \textbf{Problem 6:}
Let $C_i$ be the door with the car and $M_i$ the doors monty opens. Lets say we choose door 1, then
\[
    P(C_1) = \frac{1}{7}
\]
for staying.

We choose door 2 as the door we want to switch to, and Monty chooses 3,4,5.
\begin{align*}
    P(C_2 ~|~ M_3, M_4, M_5) & = \frac{P(M_3, M_4, M_5 ~|~ C_2) P(C_2)}{P(M_3, M_4, M_5)}                                                               \\
                             & = \frac{\frac{1}{\binom{5}{3}} * \frac{1}{7}}{\frac{1}{\binom{6}{3}}}  = \frac{\frac{1}{10} * \frac{1}{7}}{\frac{1}{20}} \\
                             & = \frac{2}{7}
\end{align*}

Yes, we should switch doors, since the probability of winning when switching is $\frac{2}{7}$ while staying is only $\frac{1}{7}$.




\end{document}